\documentclass[11pt]{article}
\usepackage[margin=1in]{geometry}
\usepackage{amsmath,amssymb,amsthm,mathtools}
\usepackage{graphicx}
\usepackage{microtype}
\usepackage[hidelinks]{hyperref}
\usepackage{xcolor}
\usepackage{enumitem}

\newtheorem{theorem}{Theorem}
\newtheorem{lemma}[theorem]{Lemma}
\newtheorem{corollary}[theorem]{Corollary}
\newtheorem{remark}[theorem]{Remark}
\newcommand{\R}{\mathbb R}
\newcommand{\Rbar}{\overline{\mathbb R}}
\newcommand{\BR}{\mathcal B_R}
\newcommand{\AR}{\mathcal A_R}
\newcommand{\diam}{\operatorname{diam}}
\newcommand{\osc}{\operatorname{osc}}

\title{\textbf{An exact finite-partition criterion for pointwise convergence into
the regulated primitive space}}
\author{Candidate full solution packet for arXiv:0911.2931}
\date{17 August 2026}

\begin{document}
\maketitle

\begin{center}
\fcolorbox{black}{gray!8}{\parbox{0.91\linewidth}{
\textbf{Status.} Full solution to the literal pointwise-convergence question,
likely valid.  The result is an intrinsic necessary-and-sufficient
characterization; bounded literature searches found no later exact answer,
but priority is not claimed.
}}
\end{center}

\section{Source question}

Erik Talvila's \emph{The regulated primitive integral}
\cite{Talvila}, arXiv:0911.2931, defines
\[
 \BR=\left\{F:\R\to\R:
 \begin{array}{l}
 F\text{ is regulated and left continuous},\\
 F(-\infty)=0,\quad F(\infty)\in\R
 \end{array}\right\}.
\]
Members are regarded as real-valued functions on
\(\Rbar=[-\infty,\infty]\), with the endpoint values equal to the respective
limits.  In the convergence section, after giving several sufficient
conditions, the paper states on PDF page~36:

\begin{center}
\includegraphics[width=0.96\linewidth]{figures/open_problem_crop.png}
\end{center}

\noindent
The central sentence asks for necessary and sufficient conditions for a
sequence in \(\BR\) to converge to a function in \(\BR\).  The surrounding
section makes clear that the convergence at issue is pointwise on
\(\Rbar\).  The theorem below gives a condition involving only the sequence,
not an already known candidate limit.

\section{Asymptotic mixed diameter}

Let \(\mathbf F=(F_n)_{n\geq1}\subset\BR\).  For a nonempty
\(E\subset\Rbar\), define its \emph{asymptotic mixed diameter}
\begin{equation}\label{eq:Delta}
 \Delta_{\mathbf F}(E)
 :=
 \sup_{x,y\in E}
 \lim_{N\to\infty}\ \sup_{n,m\geq N}
       |F_n(x)-F_m(y)|.
\end{equation}
The inner limit exists in \([0,\infty]\), since the tail suprema decrease
with \(N\).  The term mixed records that the two indices are independent.
Most importantly, the index-tail limit is taken at fixed \(x,y\)
\emph{before} the supremum over the spatial set.

An \emph{admissible partition} is a finite chain
\[
 -\infty=a_0<a_1<\cdots<a_q=\infty,
\]
whose intermediate points are real.  Its cells are
\begin{equation}\label{eq:cells}
 I_1=[-\infty,a_1],
 \qquad
 I_j=(a_{j-1},a_j]\quad(2\leq j\leq q).
\end{equation}
The right-closed convention is the one naturally adapted to left
continuity.

\section{Proof intuition}

At a fixed pair \(x,y\), pointwise convergence \(F_n\to F\) turns the
two-index tail in \eqref{eq:Delta} into \(|F(x)-F(y)|\).  Hence
\(\Delta_{\mathbf F}(E)\) becomes exactly \(\diam F(E)\).
A left-continuous regulated function is precisely a uniform limit of finite
step functions that are constant on cells of the form \eqref{eq:cells}.
This makes small mixed diameter on a finite half-open partition necessary.

For the reverse implication, a cell containing \(x\) bounds the tail
diameter of the singleton \(\{x\}\), so the condition forces
\((F_n(x))\) to be Cauchy at every point.  After taking the pointwise limit,
the same cells have small ordinary diameter.  Choosing value \(0\) on the
first cell and the limit's value at each later right endpoint builds
left-continuous regulated step functions converging uniformly to the limit.
This simultaneously enforces the two endpoint limits, all right limits, and
left continuity.

\section{The exact criterion}

We first isolate the standard half-open step characterization.

\begin{lemma}[Half-open step criterion]\label{lem:step}
Let \(F:\Rbar\to\R\) be real-valued and suppose \(F(-\infty)=0\).
Then \(F\in\BR\) if and only if, for every \(\varepsilon>0\), there is an
admissible partition whose cells satisfy
\[
 \osc(F;I_j):=\sup_{x,y\in I_j}|F(x)-F(y)|<\varepsilon
 \qquad(1\leq j\leq q).
\]
\end{lemma}

\begin{proof}
Suppose first that \(F\in\BR\).  A regulated function on the compactified
line is a uniform limit of finite step functions
\cite[\S7.13]{McLeod}; this is also the step-approximation fact used in
\cite{Talvila}.  Because \(F\) is left continuous, the values of a step
approximant at its finitely many breakpoints may be replaced by the constants
on their left.  At the first cell one uses \(0=F(-\infty)\).  Letting points
approach a breakpoint from the left shows that this modification does not
increase the limiting approximation error.  After refining at all
breakpoints, we obtain a left-continuous step function \(S\in\BR\), constant
on cells of the form \eqref{eq:cells}, with
\(\lVert F-S\rVert_\infty<\varepsilon/3\).  Thus
\[
 \osc(F;I_j)\leq 2\lVert F-S\rVert_\infty<\varepsilon.
\]

Conversely, choose such a partition with cell oscillations below
\(\varepsilon\).  Define a step function by
\[
 S_\varepsilon(x)=
 \begin{cases}
  0,&x\in I_1,\\
  F(a_j),&x\in I_j,\quad 2\leq j\leq q.
 \end{cases}
\]
Since \(-\infty\in I_1\) and \(a_j\in I_j\), the oscillation bounds give
\(\lVert F-S_\varepsilon\rVert_\infty<\varepsilon\).
The half-open convention makes \(S_\varepsilon\) left continuous; it is
regulated, equals \(0\) at the left endpoint, and is constant on a terminal
neighborhood of \(+\infty\).  Thus \(S_\varepsilon\in\BR\).
Taking \(\varepsilon\downarrow0\), uniform convergence preserves left
continuity and transfers the Cauchy criteria for right limits and endpoint
limits from \(S_\varepsilon\) to \(F\).  Therefore \(F\in\BR\).
\end{proof}

\begin{theorem}[Necessary and sufficient sequence criterion]\label{thm:main}
For a sequence \(\mathbf F=(F_n)\subset\BR\), the following are equivalent.
\begin{enumerate}[label=\textup{(\roman*)}]
\item \(F_n\) converges pointwise on \(\Rbar\) to a function \(F\in\BR\).
\item For every \(\varepsilon>0\), there is an admissible partition
      \((a_0,\dots,a_q)\) such that
      \[
       \Delta_{\mathbf F}(I_j)<\varepsilon
       \qquad(1\leq j\leq q).
      \]
\end{enumerate}
\end{theorem}

\begin{proof}
\emph{(i) implies (ii).}
Fix \(x,y\in\Rbar\).  Independent tail convergence at these two points gives
\begin{equation}\label{eq:pointpair}
 \lim_{N\to\infty}\sup_{n,m\geq N}|F_n(x)-F_m(y)|
 =|F(x)-F(y)|.
\end{equation}
Indeed, the upper bound follows from
\[
 |F_n(x)-F_m(y)|
 \leq |F_n(x)-F(x)|+|F(x)-F(y)|+|F(y)-F_m(y)|,
\]
and the reverse bound follows by applying the same inequality with the
roles of the limit and the approximants reversed.
Consequently,
\begin{equation}\label{eq:diamidentity}
 \Delta_{\mathbf F}(E)=\diam F(E)
 \quad\text{for every nonempty }E\subset\Rbar.
\end{equation}
Lemma~\ref{lem:step}, applied with the prescribed \(\varepsilon\), now gives
the required partition.

\smallskip
\noindent
\emph{(ii) implies (i).}
Fix \(x\in\Rbar\).  For every \(\varepsilon>0\), choose the partition in
(ii) and let \(I_j\) be its cell containing \(x\).  Then
\[
 \lim_{N\to\infty}\sup_{n,m\geq N}|F_n(x)-F_m(x)|
 =\Delta_{\mathbf F}(\{x\})
 \leq\Delta_{\mathbf F}(I_j)<\varepsilon.
\]
Thus \((F_n(x))_n\) is Cauchy, and the completeness of \(\R\) defines a
pointwise limit
\[
 F(x):=\lim_{n\to\infty}F_n(x)
 \qquad(x\in\Rbar).
\]
In particular \(F(-\infty)=0\), because every \(F_n(-\infty)=0\).
Identity \eqref{eq:pointpair}, and hence
\eqref{eq:diamidentity}, now holds for this \(F\).
Condition (ii) therefore gives, for every \(\varepsilon>0\), an admissible
partition on whose cells \(\osc(F;I_j)<\varepsilon\).
Lemma~\ref{lem:step} yields \(F\in\BR\), completing the proof.
\end{proof}

\begin{remark}[Metric-valued strengthening]
Nothing in the proof uses linearity of the range.  With absolute value
replaced by a metric, the theorem remains true for functions taking values
in any complete metric space, with \(F(-\infty)\) fixed at a chosen
basepoint.  In particular it holds for Banach-space-valued regulated
primitives.
\end{remark}

\section{Why the order of operations is sharp}

Define
\[
 F_n(x)=\mathbf 1_{(0,1/n]}(x),\qquad x\in\R,
 \qquad F_n(\pm\infty)=0.
\]
Each \(F_n\) is regulated and left continuous, hence belongs to \(\BR\), and
\(F_n(x)\to0\) at every \(x\in\Rbar\).  Therefore
\(\Delta_{\mathbf F}(E)=0\) for every nonempty \(E\), so
Theorem~\ref{thm:main}(ii) holds even with one cell.

This example rules out several tempting stronger replacements:
\begin{itemize}
\item \(F_n(0+)=1\) for every \(n\), so convergence of right traces is not
      necessary.
\item No tail of the sequence is equiregulated at \(0\).
\item If the spatial supremum is moved \emph{inside} the index tail,
      \[
       \widehat\Delta_{\mathbf F}(E)
       :=\lim_{N\to\infty}
       \sup_{\substack{n,m\geq N\\x,y\in E}}
          |F_n(x)-F_m(y)|,
      \]
      then every finite admissible partition has a cell containing a
      positive right-neighborhood of \(0\), and
      \(\widehat\Delta_{\mathbf F}\) equals \(1\) on that cell.  Indeed, for
      large \(n\), take \(x=1/(2n)\), \(y=2/n\), and \(m=n\).
\end{itemize}
Thus the index-limit-first ordering in \eqref{eq:Delta} is not cosmetic; it
is what admits moving or shrinking boundary layers while still detecting
the regularity of the pointwise limit.

The source paper's escaping-mass example is also detected exactly.  If
\(F_n(x)=H_1(x-n)\), then the pointwise values are \(0\) at every finite
\(x\) but \(F_n(\infty)=1\).  On the terminal cell
\((a_{q-1},\infty]\), equation \eqref{eq:Delta} compares a finite point with
\(\infty\) and gives mixed diameter at least \(1\); condition (ii) fails.

\section{Integral formulation}

\begin{corollary}\label{cor:integral}
Let \(f_n\in\AR\) have primitives \(F_n\in\BR\).  There is an
\(f=F'\in\AR\) such that
\[
 \int_{(-\infty,x)}f_n\longrightarrow
 \int_{(-\infty,x)}f
 \qquad\text{for every }x\in(-\infty,\infty]
\]
if and only if the primitive sequence satisfies
Theorem~\ref{thm:main}(ii).
\end{corollary}

\begin{proof}
Under the conventions of \cite{Talvila},
\(\int_{(-\infty,x)}f_n=F_n(x)\).  The claim is therefore exactly
Theorem~\ref{thm:main}, followed by \(f=F'\).
\end{proof}

\section{Verification, scope, and novelty}

\paragraph{Proof audit.}
The proof has only three structural inputs: completeness of the scalar
range, the fixed-pair identity \eqref{eq:pointpair}, and the elementary
half-open step characterization in Lemma~\ref{lem:step}.  The first and last
cells include the endpoint values, so the criterion enforces the endpoint
limits rather than merely local regulation.  The right endpoint of each
finite cell belongs to the cell on its left, which enforces left continuity
while allowing an arbitrary right jump.

\paragraph{Scope.}
This is a full answer to the literal question of when a sequence in
\(\BR\) converges pointwise on \(\Rbar\) to a member of \(\BR\).  It is an
intrinsic characterization rather than a conveniently checkable domination,
variation, or compactness hypothesis.  Stronger sufficient conditions in
the source remain useful precisely because they may be easier to verify.

\paragraph{Bounded novelty search.}
On 17 August 2026, the run registry, solution/attempt/proof-gap indexes, and
the local arXiv corpus were searched by arXiv id and by combinations of
the phrases pointwise convergence, regulated functions, equiregulated,
quasi-uniform, bounded epsilon-variation, and necessary and sufficient.
Web searches used the exact source title and the same close phrases.  They
found the source question, equiregulated sufficient criteria, and
approximate-variation pointwise selection results, but no later paper
stating an exact answer to this question.  Because the result is elementary
and intrinsic, an equivalent standard formulation may exist outside the
search bounds; novelty and priority are therefore not asserted.

\paragraph{Human-review recommendation.}
Verify first the quantifier/order choice in \eqref{eq:Delta}, then the
half-open endpoint conventions in Lemma~\ref{lem:step}.  The shrinking-pulse
example should be used as the quickest test that no hidden equiregulation or
trace-convergence assumption has entered.

\begin{thebibliography}{9}
\bibitem{Talvila}
E.~Talvila,
\emph{The regulated primitive integral},
Illinois J. Math. \textbf{53} (2009), no.~4, 1187--1219;
\href{https://arxiv.org/abs/0911.2931}{arXiv:0911.2931}.

\bibitem{McLeod}
R.~M. McLeod,
\emph{The Generalized Riemann Integral},
Carus Mathematical Monographs 20,
Mathematical Association of America, 1980.
\end{thebibliography}

\end{document}
